% Generated draft: requires mathematical and source review.
\documentclass[11pt]{article}
\usepackage[margin=25mm]{geometry}
\usepackage{fontspec}
\setmainfont{Latin Modern Roman}
\usepackage{amsmath,amssymb,mathtools,unicode-math}
\setmathfont{Latin Modern Math}
\usepackage{enumitem,needspace,xurl,hyperref}
\setlength{\parindent}{0pt}
\setlength{\parskip}{5pt}
\setlength{\emergencystretch}{3em}
\newcommand{\N}{\mathbb{N}}
\newcommand{\Z}{\mathbb{Z}}
\newcommand{\Q}{\mathbb{Q}}
\newcommand{\R}{\mathbb{R}}
\newcommand{\C}{\mathbb{C}}
\newcommand{\F}{\mathbb{F}}
\newcommand{\A}{\mathbb{A}}
\newcommand{\PP}{\mathbb{P}}
\newcommand{\E}{\mathbb{E}}
\newcommand{\eps}{\varepsilon}
\newcommand{\dd}{\,\mathrm{d}}
\begin{document}
\section*{UnsolvedMath ID: 30006224}
\textbf{Source problem code:} catalogue entry
\subsection{Definitions and mathematical statement}
For a smooth family of immersions \(F_t:M\to (N,g)\), let \(A_t\) denote its second fundamental form and let \(\mathbf H_t=\operatorname{tr}_{g_t}A_t\) denote its mean curvature vector, where \(g_t=F_t^*g\). Mean-curvature flow satisfies \(\partial_tF_t=\mathbf H_t\). Write
\[
B(t)=\sup_{x\in M}|\mathbf H_t(x)|.
\]
For a hypersurface, this is the supremum of the absolute value of the scalar mean curvature. A finite singular time \(T<\infty\) means a time at which the flow fails to continue smoothly in the intended solution class.

\textbf{Question as supplied, with unresolved scope made explicit.} Let \(\mathcal C_n\) stand for the geometric class of mean-curvature flows intended by the source; that class is not specified in the supplied excerpt. For every integer \(n\) with \(3<n<7\), every flow in \(\mathcal C_n\), and every finite singular time \(T\) approached by smooth flow times, must
\[
\limsup_{t\uparrow T}B(t)=+\infty?
\]
Thus the explicit dimension range in the question is \(n\in\{4,5,6\}\).

This is a schematic formulation requiring clarification of \(\mathcal C_n\), not a claim concerning all immersions into arbitrary ambient manifolds. The supplied record does not specify whether \(n\) is the dimension of the evolving manifold or its ambient space, whether the flows are hypersurface flows, or whether compactness, embeddedness, completeness, boundary conditions, or other hypotheses are intended. The displayed condition interprets ``become unbounded'' as global unboundedness near \(T\); the excerpt does not determine whether blowup is instead required locally at each singular point, or whether an actual limit of \(B(t)\) to infinity is intended.

The supplied background separately discusses dimensions \(3\) through \(6\), whereas the explicit problem uses \(3<n<7\). These ranges are retained without identifying them.
\subsection{Short English statement}
The supplied question asks whether mean curvature must become unbounded at every finite-time mean-curvature-flow singularity in the stated range \(3<n<7\). A precise theorem-level formulation requires the intended geometric hypotheses, the meaning of dimension, and the local or global meaning of blowup. The record labels the problem open; that label and its dated literature assessment are supplied curation, not independently verified conclusions.
\subsection{Sources}
\begin{itemize}
\item Supplied local UnsolvedMath registry record, permanent ID \(30006224\), source problem code \texttt{catalogue entry}, titled ``Mean-Curvature Blowup at Intermediate-Dimensional Singularities.'' The record gives proposed year \(2025\) and names no individual proposer.
\item Supplied report attribution: \emph{MATRIX-MFO Tandem Workshop: Nonlinear Geometric Diffusion Equations} (2025), Oberwolfach Reports, DOI 10.4171/owr/2025/11, pp.\ 485--530. Source URL and supplied literature reference ``Oberwolfach Report 11/2025'': \url{https://doi.org/10.4171/owr/2025/11}. The registry describes these reports as documenting mathematical workshops and their open-problem discussions; no precise page for the extracted passage is supplied.
\item Supplied literature reference: Angenent--Daskalopoulos--Sesum, ``Type-II smoothing,'' \url{https://doi.org/10.2140/apde.2026.19.857}.
\item Literature assessment checked at \texttt{2026-08-22}: the supplied curation labels the problem open and states that examples with bounded mean curvature are known in dimensions at least \(7\), while in dimensions \(3\) through \(6\) neither such examples nor an impossibility theorem are known. This is supplied dated curation derived from the OWR record, not an independently verified literature assessment. The record expressly requests independent verification of time-sensitive status and cited claims.
\item The registry assigns the statement status \texttt{corrected\_verified}, explaining that only formula delimiters, notation, grammar, and typography were repaired, with mathematical content and scope preserved. This metadata does not establish independent mathematical verification of this draft.
\end{itemize}
\end{document}
